\documentclass[border=4pt]{standalone}
\usepackage{amsmath,amssymb,mathtools,xcolor,varwidth}
\definecolor{timeblue}{HTML}{1E6FE0}
\definecolor{velgreen}{HTML}{00A35A}
\begin{document}
\begin{varwidth}{28em}
\large\bfseries
Take the very figure of Lemma IX and give it a body and a clock. Along the base we lay off      \textcolor{timeblue}{the times $AD$ and $AE$} — time runs horizontally, so a longer segment means a      longer interval since the motion began at $A$. The curve $ABC$ is now the growth of speed:      a finite force (steady, or increasing, or slackening) urges the body, so at each instant it      is moving faster than the last, exactly as Galileo saw in a falling stone. The upright      \textcolor{velgreen}{ordinates $DB$ and $EC$} are the VELOCITIES reached at times $AD$ and $AE$, drawn      at a fixed angle to the base. Thus the flat geometric picture becomes a record of accelerated      motion, with time below and speed above.
\end{varwidth}
\end{document}
